\section{Descripción del dataset}

El archivo \texttt{heart.csv} utilizado en el caso contiene 1.025 registros y 14 variables, incluida la variable objetivo \texttt{target}. No presenta valores nulos y contiene 723 filas duplicadas exactas. La clase objetivo está relativamente equilibrada, con 526 registros cuyo \texttt{target} es 1 y 499 cuyo \texttt{target} es 0. La variable \texttt{sex} presenta 713 registros con valor 1 y 312 con valor 0; la edad observada se encuentra entre 29 y 77 años.

\begin{table}[H]
  \centering
  \caption{Caracterización inicial del dataset original}
  \label{tab:caracterizacion-dataset}
  \begin{tabular}{lr}
    \toprule
    Indicador & Valor \\
    \midrule
    Registros & 1.025 \\
    Variables & 14 \\
    Valores nulos & 0 \\
    Filas duplicadas exactas & 723 (70,54\%) \\
    Registros únicos & 302 \\
    \texttt{target}=0 / \texttt{target}=1 & 499 / 526 \\
    \texttt{sex}=0 / \texttt{sex}=1 & 312 / 713 \\
    Rango de edad & 29--77 años \\
    \bottomrule
  \end{tabular}
\end{table}

Estas características orientan la factibilidad inicial, pero no validan por sí mismas la representatividad del dataset ni la pertinencia de un uso clínico. El dataset de Kaggle se utiliza como referencia académica y de prototipo; no se asume que sea la base productiva definitiva del hospital. La codificación \texttt{sex}=0 como mujer y \texttt{sex}=1 como hombre está documentada para el conjunto original de UCI; la coincidencia de valores se registra con trazabilidad externa \parencite{uci-heart-disease}.

\subsection{Diccionario de datos}

El archivo \texttt{heart.csv} contiene 14 variables: 13 predictores y la variable objetivo \texttt{target}. El diccionario documenta el significado, tipo, unidad o codificación y dominio observado de cada variable. Las definiciones conceptuales proceden de la documentación del repositorio UCI y de la versión publicada en Kaggle, mientras que dtype, rangos, nulos, valores únicos y categorías observadas se calculan directamente desde \texttt{heart.csv} \parencite{uci-heart-disease,kaggle-heart}.

Cuando la codificación original de UCI difiere del dominio presente en esta versión, se conservan los valores observados y se deja una advertencia explícita en el diccionario técnico. Esto afecta particularmente a \texttt{cp}, \texttt{slope}, \texttt{ca} y \texttt{thal}; no se asume una equivalencia que no esté documentada de forma concluyente.

\begin{table}[H]
  \centering
  \caption{Diccionario resumido de variables}
  \label{tab:diccionario-datos}
  \scriptsize
  \resizebox{\textwidth}{!}{%
  \begin{tabular}{lllll}
    \toprule
    Variable & Descripción & Tipo & Unidad / codificación & Dominio observado \\
    \midrule
    age & Edad del paciente. & Numérica discreta & años & 29--77 \\
    sex & Sexo del paciente. & Binaria & 0=mujer; 1=hombre & 0|1 \\
    cp & Tipo de dolor torácico. & Categórica & No verificada; ver CSV & 0|1|2|3 \\
    trestbps & Presión arterial en reposo al ingreso. & Numérica discreta & mmHg & 94--200 \\
    chol & Colesterol sérico. & Numérica discreta & mg/dL & 126--564 \\
    fbs & Indicador de glucosa en ayunas superior a 120 mg/dL. & Binaria & 0=no; 1=sí (>120 mg/dL) & 0|1 \\
    restecg & Resultado del electrocardiograma en reposo. & Categórica & 0=normal; 1=ST-T; 2=HVI & 0|1|2 \\
    \bottomrule
  \end{tabular}%
  }
\end{table}

\begin{table}[H]
  \centering
  \caption{Diccionario resumido de variables (continuación)}
  \label{tab:diccionario-datos-cont}
  \scriptsize
  \resizebox{\textwidth}{!}{%
  \begin{tabular}{lllll}
    \toprule
    Variable & Descripción & Tipo & Unidad / codificación & Dominio observado \\
    \midrule
    thalach & Frecuencia cardíaca máxima alcanzada. & Numérica discreta & bpm & 71--202 \\
    exang & Angina inducida por ejercicio. & Binaria & 0=no; 1=sí & 0|1 \\
    oldpeak & Depresión del segmento ST inducida por ejercicio respecto del reposo. & Numérica continua & desviación ST & 0--6.2 \\
    slope & Pendiente del segmento ST durante ejercicio. & Categórica & No verificada; ver CSV & 0|1|2 \\
    ca & Número de vasos principales coloreados por fluoroscopia. & Numérica discreta & UCI: 0--3; CSV: 0--4 & 0--4 \\
    thal & Resultado de la prueba de talio. & Categórica & No verificada; ver CSV & 0|1|2|3 \\
    target & Variable objetivo binaria de la versión Kaggle, asociada a la presencia de enfermedad cardíaca. & Binaria & Binaria 0/1; ver CSV & 0|1 \\
    \bottomrule
  \end{tabular}%
  }
\end{table}


\subsection{Alcance de los datos}

La preparación debe detectar faltantes, duplicados, valores inválidos y cambios de esquema. El análisis posterior debe evaluar calidad, estabilidad, representatividad, privacidad y posibles sesgos por subgrupos antes de que un modelo pueda considerarse para apoyo operativo.
